% called from setup-local
% \input{\pLocalSetup macros-local}

\newcommand{\urlTaylora}[0]			{http://homepage.physics.uiowa.edu/~ghowes/teach/phys7729/docs/Taylor50a.pdf}
\newcommand{\urlTaylorb}[0]			{http://homepage.physics.uiowa.edu/~ghowes/teach/phys7729/docs/Taylor50b.pdf}

\newcommand{\urlWikiTvNS}[0]		{https://en.wikipedia.org/wiki/Taylor\%E2\%80\%93von_Neumann\%E2\%80\%93Sedov_blast_wave}
\newcommand{\urlWikiTvNSSS}[0]		{\urlWikiTvNS\#Self-similar_solution}
\newcommand{\urlWikiTvNSAss}[0]		{\urlWikiTvNS\#Asymptotic_behavior_near_the_central_region}
\newcommand{\WikiTvNS}[0]			{\href{\urlWikiTvNS}{Taylor--von Neumann--Sedov Blast Wave}}

\newcommand{\urlMinpack}[0]			{https://people.sc.fsu.edu/~jburkardt/f77_src/minpack/minpack.html}

\newcommand{ \urlWikiEFE }[0]		{ https://en.wikipedia.org/wiki/Einstein_field_equations }
\newcommand{\WikiEFE}[0]			{\href{\urlWikiEFE}{Einstein Field Equations}}

\newcommand{\bpe}[0]				{\href{https://en.wikipedia.org/wiki/2020_Beirut_explosion}{Beirut Port Explosion}}

\newcommand{\curl}[1]				{\nabla \times #1}
\newcommand{\pdt}[1]				{ \pd{ #1 }{t} }
\newcommand{\pd}[2]				{ \frac{ \partial #1 } { \partial #2 } }
\newcommand{\sfrac}[0]				{ \paren{\displaystyle\tfrac{p}{\rho^{\gamma}}} }
\newcommand{\normt}[1]				{\left\lVert #1 \right\rVert_{2}}

\newcommand{\mfortran}[0]			{\href{https://www.manning.com/books/modern-fortran}{\bl{Modern Fortran}}}
\newcommand{\julia}[0]				{\href{https://julialang.org/}{\bl{Julia}}}
\newcommand{\mathematica}[0]		{\href{https://www.wolfram.com/mathematica/}{\bl{Mathematica}}}
\newcommand{\octave}[0]				{\href{https://www.gnu.org/software/octave/}{\bl{Octave}}}
\newcommand{\python}[0]				{\href{https://www.python.org/}{\bl{Python}}}

\newcommand{\mat}[2]				{\left[\begin{array} {#1}#2 \end{array}\right]}
\newcommand{\gpo}[0]				{\gamma + 1}
\newcommand{\gmo}[0]				{\gamma - 1}

% nb: /Users/dantopa/Mathematica_files/nb/projects/dtra/TvNS/taylor/paper-1-03.nb
% {11.9418, 1.01456}
% {0.0510325, 0.0188861}
\newcommand{ \solna }[0]					{ 11.942 }
\newcommand{ \solnb }[0]					{ 1.015 }
\newcommand{ \erra }[0]					{ 0.051 }
\newcommand{ \errb }[0]					{ 0.019 }
\newcommand{ \reala }[0]					{ 54474 }
\newcommand{ \realb }[0]					{ 0.3883 }

\newcommand{ \yl }[1]					{{\color{yellow} {#1}}}
\newcommand{ \blk }[1]					{{\color{black} {#1}}}

\DeclareMathOperator*{\argmin}{arg\,min}

\endinput  %  ==  ==  ==  ==  ==  ==  ==  ==  ==

\newcommand{  }[0]					{  }
